\section{Presence, Vection, and Stationary Self-Motion}
\label{sec:background_presence_vection_stationary_motion}

Embodiment is related to, but not identical with, presence. Presence concerns the experience of being in a mediated environment, while embodiment concerns the experienced relation between the self and a body or body part. In VR, these experiences often interact. A believable environment may make the virtual body feel more situated, and a coherent body may make the environment feel more actionable. Nevertheless, the concepts should remain separate because a user can feel present in a scene without strongly experiencing a virtual body as their own.
% TODO cite: presence survey / Skarbez et al.
% TODO cite: Slater place illusion and plausibility illusion.

This distinction is important for VR rowing because the boat-centred mode may increase the plausibility of the sport scenario. Seeing water, a shell, oars, and forward movement may make the situation feel more like rowing. This does not automatically mean that embodiment is stronger. A scene can be convincing as a rowing environment while the virtual hands or body still feel less directly connected to the user's own movement.

The visual experience of moving through space is also relevant. In stationary exercise systems, the body remains in one physical location while the virtual environment may depict locomotion. Optic flow can create a visual impression of self-motion, and vection describes the illusion of self-motion in a stationary observer. In rowing, this can make an ergometer action feel connected to movement through a river or race course, even though the user remains in the room.
% TODO cite: optic flow definition/history.
% TODO cite: Palmisano et al. or Kooijman et al. for vection.

For this thesis, vection is not the main dependent construct, but it helps explain why a boat-centred representation may feel meaningful. The user's effort produces visible locomotion. The stroke is therefore not only an isolated body movement, but appears to move the user through a virtual environment. This can support plausibility and motivation, and may also influence self-location by encouraging the user to experience themselves as situated in the virtual rowing scene.

At the same time, stationary self-motion is a design tension. Visual motion that is not sufficiently related to the user's action can create a mismatch between visual and bodily cues. In many VR applications, such mismatches are discussed in relation to cybersickness. In rowing, the risk may be moderated by the fact that visible motion is coupled to the user's own rhythmic effort rather than imposed independently. However, this should be treated cautiously and empirically rather than assumed.
% TODO cite: VR locomotion/cybersickness source if retained.
% TODO cite: vection and visually induced self-motion literature if used.

The two rowing representations therefore differ not only in body mapping, but also in how strongly they situate the user in a moving sport scenario. The ergometer-congruent mode may offer stronger correspondence with the physical body and machine. The boat-centred mode may offer stronger environmental plausibility and visually supported self-motion. This makes presence and vection useful background concepts, while the thesis remains focused on embodiment and its subcomponents.
